\documentclass[10pt,a4paper]{amsart}
\usepackage{amsmath,amssymb,amsthm,graphicx}
% Retained preamble for the Stacks Project chapter divisors.tex (Divisors).
% The chapter file pinned-original-source.tex carries only \input{preamble} on its first line, so no
% preamble text lives inside the chapter itself.  The Stacks Project uses one single shared preamble for
% every chapter; the snapshot stages it verbatim in sources/stacks-chow/preamble.tex.  The block below is
% copied mechanically, line for line, from that file's lines 14-18 (the Xy-pic setup) and lines 154-217
% (the native theorem declarations together with the mathematical macros).  Nothing was added, removed or
% redefined.  The wrapper class and the amsmath/amssymb/amsthm/graphicx base packages are supplied by the
% assembler, so the class/\externaldocument part of the shared preamble is not repeated here.
% For commutative diagrams we use Xy-pic
\usepackage[all]{xy}

% We use 2cell for 2-commutative diagrams.
\xyoption{2cell}
% Theorem environments.
%
\theoremstyle{plain}
\newtheorem{theorem}[subsection]{Theorem}
\newtheorem{proposition}[subsection]{Proposition}
\newtheorem{lemma}[subsection]{Lemma}

\theoremstyle{definition}
\newtheorem{definition}[subsection]{Definition}
\newtheorem{example}[subsection]{Example}
\newtheorem{exercise}[subsection]{Exercise}
\newtheorem{situation}[subsection]{Situation}

\theoremstyle{remark}
\newtheorem{remark}[subsection]{Remark}
\newtheorem{remarks}[subsection]{Remarks}

\numberwithin{equation}{subsection}

% Macros
%
\def\lim{\mathop{\mathrm{lim}}\nolimits}
\def\colim{\mathop{\mathrm{colim}}\nolimits}
\def\Spec{\mathop{\mathrm{Spec}}}
\def\Hom{\mathop{\mathrm{Hom}}\nolimits}
\def\Ext{\mathop{\mathrm{Ext}}\nolimits}
\def\SheafHom{\mathop{\mathcal{H}\!\mathit{om}}\nolimits}
\def\SheafExt{\mathop{\mathcal{E}\!\mathit{xt}}\nolimits}
\def\Sch{\mathit{Sch}}
\def\Mor{\mathop{\mathrm{Mor}}\nolimits}
\def\Ob{\mathop{\mathrm{Ob}}\nolimits}
\def\Sh{\mathop{\mathit{Sh}}\nolimits}
\def\NL{\mathop{N\!L}\nolimits}
\def\CH{\mathop{\mathrm{CH}}\nolimits}
\def\proetale{{pro\text{-}\acute{e}tale}}
\def\etale{{\acute{e}tale}}
\def\QCoh{\mathit{QCoh}}
\def\Ker{\mathop{\mathrm{Ker}}}
\def\Im{\mathop{\mathrm{Im}}}
\def\Coker{\mathop{\mathrm{Coker}}}
\def\Coim{\mathop{\mathrm{Coim}}}

% Boxtimes
%
\DeclareMathSymbol{\boxtimes}{\mathbin}{AMSa}{"02}

%
% Macros for moduli stacks/spaces
%
\def\QCohstack{\mathcal{QC}\!\mathit{oh}}
\def\Cohstack{\mathcal{C}\!\mathit{oh}}
\def\Spacesstack{\mathcal{S}\!\mathit{paces}}
\def\Quotfunctor{\mathrm{Quot}}
\def\Hilbfunctor{\mathrm{Hilb}}
\def\Curvesstack{\mathcal{C}\!\mathit{urves}}
\def\Polarizedstack{\mathcal{P}\!\mathit{olarized}}
\def\Complexesstack{\mathcal{C}\!\mathit{omplexes}}
% \Pic is the operator that assigns to X its picard group, usage \Pic(X)
% \Picardstack_{X/B} denotes the Picard stack of X over B
% \Picardfunctor_{X/B} denotes the Picard functor of X over B
\def\Pic{\mathop{\mathrm{Pic}}\nolimits}
\def\Picardstack{\mathcal{P}\!\mathit{ic}}
\def\Picardfunctor{\mathrm{Pic}}
\def\Deformationcategory{\mathcal{D}\!\mathit{ef}}

% --- appended 2026-10-12 (Chief): the project own definitions of the reference/slogan/history
% display environments, copied mechanically from the shared Stacks preamble lines 8-12
% (\usepackage{verbatim} plus the three \newenvironment lines, which the chapter preamble keeps
% verbatim from the project preamble). Nothing added, removed or redefined: these environments
% carry no mathematical content and the project itself typesets them as comments. Without these
% three lines a counted statement whose own text uses them fails with Environment slogan undefined.
\usepackage{verbatim}
\newenvironment{reference}{\comment}{\endcomment}
\newenvironment{slogan}{\comment}{\endcomment}
\newenvironment{history}{\comment}{\endcomment}
\usepackage{cleveref}
\crefname{originallemma}{original lemma}{original lemmas}
\Crefname{originallemma}{Original lemma}{Original lemmas}
\makeatletter
\newlabel{algebra-lemma-total-ring-fractions-no-embedded-points}{{02LX}{source}{}{sourcepointer-1}{}}
\newlabel{algebra-lemma-total-ring-fractions-no-embedded-points@cref}{{[originallemma][0][]02LX}{[1][0][]source}{}{}{}}
\makeatother
\crefname{originallemma}{original lemma}{original lemmas}
\Crefname{originallemma}{Original lemma}{Original lemmas}
\makeatletter
\newlabel{algebra-lemma-weakly-ass-zero-divisors}{{05C3}{source}{}{sourcepointer-2}{}}
\newlabel{algebra-lemma-weakly-ass-zero-divisors@cref}{{[originallemma][0][]05C3}{[1][0][]source}{}{}{}}
\makeatother
\crefname{originallemma}{original lemma}{original lemmas}
\Crefname{originallemma}{Original lemma}{Original lemmas}
\makeatletter
\newlabel{lemma-pullback-meromorphic-sections-defined}{{02OU}{source}{}{sourcepointer-3}{}}
\newlabel{lemma-pullback-meromorphic-sections-defined@cref}{{[originallemma][0][]02OU}{[1][0][]source}{}{}{}}
\makeatother
\crefname{originaldefinition}{original definition}{original definitions}
\Crefname{originaldefinition}{Original definition}{Original definitions}
\makeatletter
\newlabel{morphisms-definition-rational-function}{{01RT}{source}{}{sourcepointer-4}{}}
\newlabel{morphisms-definition-rational-function@cref}{{[originaldefinition][0][]01RT}{[1][0][]source}{}{}{}}
\makeatother
\crefname{originallemma}{original lemma}{original lemmas}
\Crefname{originallemma}{Original lemma}{Original lemmas}
\makeatletter
\newlabel{morphisms-lemma-integral-scheme-rational-functions}{{01RV}{source}{}{sourcepointer-5}{}}
\newlabel{morphisms-lemma-integral-scheme-rational-functions@cref}{{[originallemma][0][]01RV}{[1][0][]source}{}{}{}}
\makeatother
\crefname{originallemma}{original lemma}{original lemmas}
\Crefname{originallemma}{Original lemma}{Original lemmas}
\makeatletter
\newlabel{schemes-lemma-push-forward-quasi-coherent}{{01LC}{source}{}{sourcepointer-6}{}}
\newlabel{schemes-lemma-push-forward-quasi-coherent@cref}{{[originallemma][0][]01LC}{[1][0][]source}{}{}{}}
\makeatother
\crefname{originalsection}{original section}{original sections}
\Crefname{originalsection}{Original section}{Original sections}
\makeatletter
\newlabel{schemes-section-points}{{00Y3}{source}{}{sourcepointer-7}{}}
\newlabel{schemes-section-points@cref}{{[originalsection][0][]00Y3}{[1][0][]source}{}{}{}}
\makeatother
\title{Under weak association and locally finite components the sheaf of meromorphic functions is the product of the skyscrapers at the generic points}
\author{The Stacks Project Authors}
\date{stacks-project:divisors.tex:local-snapshot-20261012; GNU FDL 1.2 or later}
\begin{document}
\maketitle
\noindent\textit{Editorial note.} All delivered mathematical text is the author's own, copied from the pinned Stacks Project chapter file pinned-original-source.tex; no statement or proof step was rewritten, reordered or removed.\par\medskip
\noindent\textit{Editorial note.} Following the multi-claim convention for this source, source\_url is the source-level URL https://stacks.math.columbia.edu/ and the source label lemma-meromorphic-weakass-finite is carried in claim\_key and here rather than in a per-claim tag URL.\par\medskip
\noindent\textit{Editorial note.} The preamble is the chapter's shared Stacks Project preamble, retained as sources/stacks-divisors/retained-preamble.tex: a line-for-line excerpt of the shared Stacks preamble staged with this snapshot. No macro was invented or redefined, and this card adds no preamble text of its own.\par\medskip
\noindent\textit{Editorial note.} Statement span 6034-6073 and proof span 6075-6132 come from the halfA pool card operations/pool-cards/stacks-divisors-halfA.json and were re-read in the pinned chapter file before this card was written.\par\medskip
\noindent\textit{Editorial note.} The proof contains the author's parenthesis '(argument omitted)' for the extension of the rational-function comparison from the case of finitely many generic points to the general case by glueing; that phrase is kept verbatim and the identity it depends on is proved in the counted span.\par\medskip
\noindent\textit{Editorial note.} The statement carries a long list of numbered consequences; all of them are part of the counted statement span and are preserved verbatim.\par\medskip
\noindent\textit{Editorial note.} 7 cross-references occur inside the counted spans; they are listed in external\_references as navigation pointers only.\par\medskip
\noindent\textit{Allowed standard prerequisites (editorial).}\begin{itemize}
\item The description of the total ring of fractions of a ring all of whose weakly associated primes are minimal, and the identification of nonzerodivisors with elements outside the weakly associated primes. Cited at proof lines 6102-6106 and again at 6126-6129. Gives the product decomposition of the total ring of fractions over affine opens and the statement about nonzerodivisors.
\item The compatibility of the sheaf of meromorphic functions with tensor products by quasi-coherent modules, and the identification of rational functions with meromorphic functions on an integral scheme. Used at proof lines 6077-6082. Derives the statement about meromorphic sections of an arbitrary quasi-coherent module from the main identity, and handles part of the comparison with rational functions.
\item Pullbacks of meromorphic functions along a morphism whose fibres are integral, cited from the same chapter, together with quasi-compactness and quasi-separatedness of the coproduct morphism and quasi-coherence of its pushforward of the structure sheaf. Cited at proof lines 6088-6090 and 6120-6123. Produces the comparison map from the sheaf of meromorphic functions to the pushforward and gives its quasi-coherence.
\end{itemize}\medskip
\ifdefined\hypertarget\hypertarget{sourcepointer-1}{}\fi
\noindent\textit{Source reference.} Original lemma 02LX: focused-seeds/AGP/sources/stacks-snapshot-20261011/tags.txt, source lines 3424--3424. Same-chapter reference used in the retained author text at source line(s) 6105. It is supplied only as a navigation pointer so that the retained text compiles; the counted item delivers the author's own statement and the author's complete proof, and no part of the referenced item is copied in or substituted for the counted argument.\par\medskip
\ifdefined\hypertarget\hypertarget{sourcepointer-2}{}\fi
\noindent\textit{Source reference.} Original lemma 05C3: focused-seeds/AGP/sources/stacks-snapshot-20261011/tags.txt, source lines 6906--6906. Same-chapter reference used in the retained author text at source line(s) 2624, 6025, 6106, 6127. It is supplied only as a navigation pointer so that the retained text compiles; the counted item delivers the author's own statement and the author's complete proof, and no part of the referenced item is copied in or substituted for the counted argument.\par\medskip
\ifdefined\hypertarget\hypertarget{sourcepointer-3}{}\fi
\noindent\textit{Source reference.} Original lemma 02OU: focused-seeds/AGP/sources/stacks-snapshot-20261011/tags.txt, source lines 3529--3529. Same-chapter reference used in the retained author text at source line(s) 6096, 6992. It is supplied only as a navigation pointer so that the retained text compiles; the counted item delivers the author's own statement and the author's complete proof, and no part of the referenced item is copied in or substituted for the counted argument.\par\medskip
\ifdefined\hypertarget\hypertarget{sourcepointer-4}{}\fi
\noindent\textit{Source reference.} Original definition 01RT: focused-seeds/AGP/sources/stacks-snapshot-20261011/tags.txt, source lines 2314--2314. Same-chapter reference used in the retained author text at source line(s) 6069. It is supplied only as a navigation pointer so that the retained text compiles; the counted item delivers the author's own statement and the author's complete proof, and no part of the referenced item is copied in or substituted for the counted argument.\par\medskip
\ifdefined\hypertarget\hypertarget{sourcepointer-5}{}\fi
\noindent\textit{Source reference.} Original lemma 01RV: focused-seeds/AGP/sources/stacks-snapshot-20261011/tags.txt, source lines 2316--2316. Same-chapter reference used in the retained author text at source line(s) 1879, 6083. It is supplied only as a navigation pointer so that the retained text compiles; the counted item delivers the author's own statement and the author's complete proof, and no part of the referenced item is copied in or substituted for the counted argument.\par\medskip
\ifdefined\hypertarget\hypertarget{sourcepointer-6}{}\fi
\noindent\textit{Source reference.} Original lemma 01LC: focused-seeds/AGP/sources/stacks-snapshot-20261011/tags.txt, source lines 2081--2081. Same-chapter reference used in the retained author text at source line(s) 741, 6120, 7924. It is supplied only as a navigation pointer so that the retained text compiles; the counted item delivers the author's own statement and the author's complete proof, and no part of the referenced item is copied in or substituted for the counted argument.\par\medskip
\ifdefined\hypertarget\hypertarget{sourcepointer-7}{}\fi
\noindent\textit{Source reference.} Original section 00Y3: focused-seeds/AGP/sources/stacks-snapshot-20261011/tags.txt, source lines 1241--1241. Same-chapter reference used in the retained author text at source line(s) 6050, 6325, 6435. It is supplied only as a navigation pointer so that the retained text compiles; the counted item delivers the author's own statement and the author's complete proof, and no part of the referenced item is copied in or substituted for the counted argument.\par\medskip
\begin{lemma}
\label{lemma-meromorphic-weakass-finite}
Let $X$ be a scheme such that
\begin{enumerate}
\item[(a)] every weakly associated point of $X$ is a generic point of an
irreducible component of $X$, and
\item[(b)] any quasi-compact open has a finite number of irreducible components.
\end{enumerate}
Let $X^0$ be the set of generic points of irreducible components of $X$.
Then we have
$$
\mathcal{K}_X =
\bigoplus\nolimits_{\eta \in X^0} j_{\eta, *}\mathcal{O}_{X, \eta} =
\prod\nolimits_{\eta \in X^0} j_{\eta, *}\mathcal{O}_{X, \eta}
$$
where $j_\eta : \Spec(\mathcal{O}_{X, \eta}) \to X$ is the canonical map
of Schemes, Section \ref{schemes-section-points}. Moreover
\begin{enumerate}
\item $\mathcal{K}_X$ is a quasi-coherent sheaf of
$\mathcal{O}_X$-algebras,
\item for every quasi-coherent $\mathcal{O}_X$-module $\mathcal{F}$ the sheaf
$$
\mathcal{K}_X(\mathcal{F}) =
\bigoplus\nolimits_{\eta \in X^0} j_{\eta, *}\mathcal{F}_\eta =
\prod\nolimits_{\eta \in X^0} j_{\eta, *}\mathcal{F}_\eta
$$
of meromorphic sections of $\mathcal{F}$
is quasi-coherent,
\item $\mathcal{S}_x \subset \mathcal{O}_{X, x}$
is the set of nonzerodivisors for any $x \in X$,
\item $\mathcal{K}_{X, x}$ is the total quotient ring of $\mathcal{O}_{X, x}$
for any $x \in X$,
\item $\mathcal{K}_X(U)$ equals the total quotient ring of $\mathcal{O}_X(U)$
for any affine open $U \subset X$,
\item the ring of rational functions of $X$
(Morphisms, Definition \ref{morphisms-definition-rational-function})
is the ring of meromorphic
functions on $X$, in a formula: $R(X) = \Gamma(X, \mathcal{K}_X)$.
\end{enumerate}
\end{lemma}

\begin{proof}
Observe that a locally finite direct sum of sheaves of modules
is equal to the product since you can check this on stalks for
example. Then since $\mathcal{K}_X(\mathcal{F}) =
\mathcal{F} \otimes_{\mathcal{O}_X} \mathcal{K}_X$
we see that (2) follows from the other statements.
Also, observe that part (6) follows from the initial
statement of the lemma and Morphisms, Lemma
\ref{morphisms-lemma-integral-scheme-rational-functions}
when $X^0$ is finite; the general case of (6) follows from this
by glueing (argument omitted).

\medskip\noindent
Let $j : Y = \coprod\nolimits_{\eta \in X^0} \Spec(\mathcal{O}_{X, \eta}) \to X$
be the coproduct of the morphisms $j_\eta$. We have to show that
$\mathcal{K}_X = j_*\mathcal{O}_Y$.
First note that $\mathcal{K}_Y = \mathcal{O}_Y$ as $Y$ is a disjoint
union of spectra of local rings of dimension $0$: in a local
ring of dimension zero any nonzerodivisor is a unit.
Next, note that pullbacks of meromorphic
functions are defined for $j$ by
Lemma \ref{lemma-pullback-meromorphic-sections-defined}.
This gives a map
$$
\mathcal{K}_X \longrightarrow j_*\mathcal{O}_Y.
$$
Let $\Spec(A) = U \subset X$ be an affine open. Then $A$ is a ring
with finitely many minimal primes $\mathfrak q_1, \ldots, \mathfrak q_t$
and every weakly associated prime of $A$ is one of the $\mathfrak q_i$.
We obtain $Q(A) = \prod A_{\mathfrak q_i}$
by Algebra, Lemmas \ref{algebra-lemma-total-ring-fractions-no-embedded-points}
and \ref{algebra-lemma-weakly-ass-zero-divisors}.
In other words, already the value of the presheaf
$U \mapsto \mathcal{S}(U)^{-1}\mathcal{O}_X(U)$ agrees with
$j_*\mathcal{O}_Y(U)$ on our affine open $U$. Hence the displayed
map is an isomorphism which proves the first displayed equality in
the statement of the lemma.

\medskip\noindent
Finally, we prove (1), (3), (4), and (5).
Part (5) we saw during the course of the proof that
$\mathcal{K}_X = j_*\mathcal{O}_Y$.
The morphism $j$ is quasi-compact by our assumption
that the set of irreducible components of $X$ is locally finite.
Hence $j$ is quasi-compact and quasi-separated (as $Y$ is separated).
By Schemes, Lemma \ref{schemes-lemma-push-forward-quasi-coherent}
$j_*\mathcal{O}_Y$ is quasi-coherent. This proves (1).
Let $x \in X$. We may choose an affine open neighbourhood
$U = \Spec(A)$ of $x$ all of whose irreducible components
pass through $x$. Then $A \subset A_\mathfrak p$ because every
weakly associated prime of $A$ is contained in $\mathfrak p$
hence elements of $A \setminus \mathfrak p$ are nonzerodivisors
by Algebra, Lemma \ref{algebra-lemma-weakly-ass-zero-divisors}.
It follows easily that any nonzerodivisor of $A_\mathfrak p$
is the image of a nonzerodivisor on a (possibly smaller)
affine open neighbourhood of $x$. This proves (3).
Part (4) follows from part (3) by computing stalks.
\end{proof}

\end{document}
